\documentclass[11pt]{amsart}
\usepackage[T1]{fontenc}
\usepackage{lmodern,amsmath,amssymb,mathtools,mathrsfs,booktabs,microtype}
\usepackage[margin=1.08in]{geometry}
\usepackage[hidelinks]{hyperref}
\pdfinfoomitdate=1\pdftrailerid{}\pdfsuppressptexinfo=15
\hypersetup{pdftitle={General Theta Foundations I: Acquired Geometry and Causal Resource Transfer},pdfauthor={Qian Qi}}
\newtheorem{theorem}{Theorem}[section]
\newtheorem{proposition}[theorem]{Proposition}
\newtheorem{lemma}[theorem]{Lemma}
\newtheorem{corollary}[theorem]{Corollary}
\theoremstyle{definition}\newtheorem{definition}[theorem]{Definition}
\theoremstyle{remark}\newtheorem{remark}[theorem]{Remark}
\newcommand{\E}{\mathbb E}\newcommand{\PP}{\mathbb P}\newcommand{\R}{\mathbb R}
\newcommand{\law}{\operatorname{Law}}\newcommand{\TV}{\operatorname{TV}}
\newcommand{\Lip}{\operatorname{Lip}}\newcommand{\tr}{\operatorname{tr}}
\newcommand{\supp}{\operatorname{supp}}\newcommand{\aff}{\operatorname{Aff}}
\newcommand{\cP}{\mathcal P}\newcommand{\cx}{\preceq_{\rm cx}}
\newcommand{\ext}{\mathrm{ext}}\newcommand{\aut}{\mathrm{aut}}\newcommand{\cp}{\mathrm{cp}}
\newcommand{\norm}[1]{\left\lVert#1\right\rVert}
\title{General Theta Foundations I:\\Acquired Geometry and Causal Resource Transfer}
\author{Qian Qi}\date{9 October 2026}
\subjclass[2020]{Primary 62B05, 60G35; Secondary 60E15, 93E20, 94A34, 81P45}
\keywords{causal experiment, predictive quotient, convex order, finite-state observer, acquired geometry, quantum instrument}
\begin{document}
\begin{abstract}
Which distributions of predictive states can a finite causal observer actually retain? We give an exact allocation criterion on compact metrizable predictive carriers derived from prepared experimental kernels. With an external stage index, realizability is equivalent to successive finite-support barycentric contractions of the deployed acquisition laws. For a single autonomous machine, separate contractions are insufficient: a simultaneous affine-test inequality characterizes when all phases admit the same report-to-label kernels. The resulting variational formulas equal the optimal decision risks. Every infeasible proposed allocation has a finite affine dual witness, and compactness gives global finite obstructions for external width profiles. An exact action-slack and Jensen decomposition identifies zero-loss completion without density, dimension or contraction assumptions. Quantitative consequences recover matching memory laws for uniformly smoothing, block-stable experiments, with distinct external-clock and charged-clock accounts. A noncontracting capture-and-blind-transport experiment has instead an exact risk governed by the smallest intermediate alphabet; fractal and countably atomic acquisition laws give polynomial and exponential examples. Gaussian hidden-state observations and noncommuting finite-dimensional instruments verify the regular hypotheses for ordinary probability prediction. Numerical and causal-simulation errors are composed only with explicit execution and common-task certificates.
\end{abstract}
\maketitle
\section{The completion problem}\label{sec:intro}
A terminal statistic may be inexpensive to store and impossible to maintain with the same online information. The missing object is not a dimension of the formal state space, but a compatibility condition between acquisition and every intermediate act of forgetting. This paper studies that condition for a prepared causal experiment and a specified decision task.

Our starting data are positive normalized kernels, legal interventions, and physically executable future tests. Their expectations identify histories that have the same predictive content. Mixtures of preparations give a convex predictive carrier, whereas an implemented acquisition policy gives a probability law on that carrier. These two objects are different. Compression replaces an acquired state by a conditional barycenter. It therefore contracts the acquired law in convex order, but a sequence of such contractions does not by itself describe an autonomous finite machine: the same persistent label cannot invoke different transition programmes at different unobserved times.

Theorem~\ref{thm:completion} gives an exact answer to both issues. For an externally clocked observer, a finite-support law and an action distribution at each retained state suffice; one tests a one-step barycentric inequality at each cut. For an autonomous observer, one keeps the proposed label occupancies and their conditional predictive states at every cut. A simultaneous dual inequality tests whether a single label/action/report transition rule can realize all of them. It eliminates the unknown transition kernels from the feasibility condition, without presuming a phase index or an exact posterior register. The criterion is variational, not a closed-form exponent or an efficient synthesis algorithm. A local converse certifies an infeasible proposed law sequence. Theorem~\ref{thm:finite_obstruction} additionally gives optimal external observers and finite test families excluding every proposed flow below an unattainable risk threshold. This global compactness conclusion has no automatic effective counterpart for general stationary Borel programmes.

The exact mechanism uses no acquired density, finite dimension, smooth task, summable stability, or common report support. Theorem~\ref{thm:completion} gives the risk frontier and Corollary~\ref{cor:zero} gives an exact zero-loss criterion. In contrast, a power-law conclusion needs quantitative hypotheses. Under an all-state/all-action density bound, block summability, and strong curvature of the task, Theorem~\ref{thm:regular} gives, uniformly in the horizon $n$,
\begin{equation}\label{eq:intro_regular}
 R_{\cp}(n,M)-B_n^*\asymp R_{\ext}(n,M)-B_n^*\asymp M^{-2/d},
 \qquad R_{\aut}(n,M)-B_n^*\asymp(M-n)^{-2/d}.
\end{equation}
The last formula assumes clock-neutral report supports and an internally exact deadline; it is not the conventional externally clocked filtering law. The all-state density condition is stronger than positive observation noise, and the strong-curvature condition is absent from the exact criterion.

A different sharp law makes the distinction substantive. When one capture reveals a prepared variable and all later reports are blind, Theorem~\ref{thm:blind} proves
\begin{equation}\label{eq:intro_blind}
 R_{\ext}(\mathbf m)=q_{\min_t m_t}(\mu),\qquad
 R_{\aut}(n,M)=q_{\lfloor(M-1)/n\rfloor}(\mu).
\end{equation}
Here $q_L(\mu)$ is the squared quantization error of the \emph{actually captured} distribution. The predictor must carry its resolution through every intervening cut; spending almost all labels at the terminal cut cannot restore discarded information. The acquired law may be fractal or countably atomic. This noncontracting class exhibits a different memory law without leaving the same general criterion.

The Gaussian and quantum realizations in Section~\ref{sec:realizations} establish the regular hypotheses from raw kernels for familiar probability forecasts: categorical filtering under Brier loss and prediction of a fixed informationally complete measurement. The latter is a finite-dimensional, known-instrument theorem; noncommutativity does not imply a result for arbitrary quantum control. The numerical statements give error propagation \emph{given} integration, selection and execution certificates, not polynomial-time synthesis or universal computability.

The one-step convex-order theorem is classical in substance\cite{strassen,blackwell}; conditional means, Bellman recursion and convex separation are not claimed as new. The contribution proposed here is their exact causal composition with counted finite supports, simultaneous phase-reuse constraints, an action-slack identity under the deployed law, and matching interface-dependent consequences. The mass--curvature and block-stable arguments are inherited and re-proved in the form used here\cite{prior}. Complete earlier articles are retained as archival companions, but no proof in this article requires a realization theorem or an unproved assertion from them.

\section{Experimental data and observers}\label{sec:experiment}
We use $\TV(P,Q)=\sup_B|P(B)-Q(B)|=\frac12\|P-Q\|_1$. Fix a finite horizon $n\ge1$. A call takes a hidden physical state $x$, a legal action $a$, and produces a new hidden state, a report, and a resource increment through a positive normalized kernel. Positive means nonnegative and normalized; deterministic reports and zero likelihoods are permitted. Preparation is fixed. An unknown parameter, when present, is not an argument available to the observer. The risk results fix a known prepared law; adjoining an unknown parameter to the hidden state gives a Bayes experiment, not an automatic frequentist minimax theorem.

An executable future test is a legal continuation with a bounded scored outcome. At a fixed phase, histories are equivalent if all such tests have the same conditional expectation and the same legal interface. Appendix~\ref{sec:quotient} gives sufficient raw-kernel hypotheses for a compact metrizable realization of this quotient. In particular it establishes measurability, actual mixture realizability, and minimality. We now state precisely the properties used by the main theorem.

For each $t=0,\ldots,n$, let $K_t$ be a compact metrizable convex subset of a Hausdorff locally convex space, embedded by a countable separating family of continuous affine future-test expectations. Write $\aff(K_t)$ for continuous affine functions, including constants. Probability measures on $K_t$ have barycenters in $K_t$, denoted $\bar\nu$, and $f(\bar\nu)=\int f\,d\nu$ for $f\in\aff(K_t)$. No finite dimension is assumed. Let $\mathcal A$ be finite and let $\mathcal A_t\subseteq\mathcal A$ be a fixed nonempty legal action set at phase $t<n$. Let $\mathcal Y_a$ be standard Borel. These sets are part of the physical interface, not a phase variable readable by a stationary machine. Kernels for illegal actions may be extended arbitrarily and never enter a positive-mass term. For $t<n$, the raw instruments induce jointly Borel maps
\begin{equation}\label{eq:instrument}
 J_t(q,a)(dy),\qquad F_t(q,a,y)\in K_{t+1},\qquad
 \Lambda_t(q,a)=\law_{Y\sim J_t(q,a)}F_t(q,a,Y).
\end{equation}
Values of $F_t$ on zero-likelihood reports are irrelevant and are assigned measurably. The instrument property means that, for every probability $\nu$ on $K_t$, Borel report set $B$, and $f\in\aff(K_{t+1})$,
\begin{align}\label{eq:affine_instrument}
 J_t(\bar\nu,a)(B)&=\int J_t(q,a)(B)\nu(dq),\\
 \int_B f(F_t(\bar\nu,a,y))J_t(\bar\nu,a)(dy)
 &=\int\!\int_B f(F_t(q,a,y))J_t(q,a)(dy)\nu(dq).\notag
\end{align}
The identities are required on all mixture states in the carrier. They express physical mixing before an unnormalized instrument, not linearity of the normalized Bayes update. Assume $q\mapsto\Lambda_t(q,a)$ is weakly continuous. This ensures continuous Bellman values; the feasibility parts below only need the displayed measurable instruments and conditional barycenters. Total-variation report continuity is an additional quantitative property verified in the regular realizations, not a premise of convex order.

The terminal decisions form a compact metric space $U$. The scored loss has conditional expectation $\ell(q,u)\in[0,H]$, jointly continuous and affine in $q$ for each $u$. Let
\begin{equation}\label{eq:bellman}
 g(q)=\min_{u\in U}\ell(q,u),\quad V_n=g,\quad
 V_t(q)=\min_{a\in\mathcal A_t}\int V_{t+1}(w)\Lambda_t(q,a)(dw),\quad B_n^*=V_0(q_0).
\end{equation}
Compactness, weak continuity and finite actions give continuous values and Borel minimizing decisions/actions. Concavity follows by taking infima of the affine risks of fixed continuation programmes: a programme may depend on its subsequently observed reports, but is fixed before the incoming preparation is mixed. Its affine risk need not be continuous for this argument. The instrument identities make that mixture comparison valid. The terminal task and hence $B_n^*$ may change with the horizon.

\begin{definition}[Two finite-state interfaces]\label{def:machines}
An externally clocked observer has finite sets $Z_t$, with $|Z_0|=1$ and $|Z_t|\le m_t$ for $1\le t\le n$. It uses action probabilities $\alpha_t(a\mid z)$ and a Borel stochastic kernel $k_t(z'\mid z,a,y)$. The external phase $t$ is free in this model. Terminal decisions depend only on $Z_n$.

An autonomous observer is a \emph{single} tuple $(Z,z_0,D,\alpha,k,u)$: $|Z|\le M$, $D\subset Z$ is the set of stopping labels, $\alpha(a\mid z)$ is one time-independent action kernel on $Z\setminus D$, $k(z'\mid z,a,y)$ is one time-independent next-label kernel, and $u:D\to U$ is one terminal readout map. At a stopping label it issues the readout and makes no further call. Starting from $z_0$, it must stop after exactly $n$ calls almost surely. Neither kernel may consult absolute time unless it is encoded in $z$. Fresh independent randomization is erased at the next persistent cut. A retained pseudorandom-generator seed, instruction pointer, output tape or clock is part of $Z$.
\end{definition}
A fixed read-only programme can depend on the known preparation, kernels, horizon and budgets, but not on the acquired data or an inaccessible parameter. Programme length, temporary workspace, precision, physical calls and physical time are separate resources. Exact-real measurable programmes define the cardinality theorem; finite-bit implementation is addressed in Section~\ref{sec:resources}. An unused label may have arbitrary programmed behavior. Legality means $\alpha_t(\mathcal A_t\mid z)=1$ on every positive-mass phase label; for a stationary observer the same $\alpha_z$ must obey this requirement at every phase at which $z$ occurs. All measures below are generated by the actual deployed observer, including its fresh randomization.

Write $R_{\ext}(\mathbf m)$ and $R_{\aut}(n,M)$ for the infimal risks in these two classes, with infimum $+\infty$ for an empty class. Write $R_{\ext}(n,M)$ when every $m_t=M$. The checkpoint relaxation permits full history during acquisition and restricts only the final label to $M$ values. These are different information patterns, all using the same preparation and task.

\section{A causal allocation theorem}\label{sec:completion}
For probabilities $\nu,\lambda$ on a compact convex carrier, write $\nu\cx\lambda$ when
$\int f\,d\nu\le\int f\,d\lambda$ for every continuous convex $f$. Let $\mathcal P_m(K)$ be the probabilities with at most $m$ atoms, not the measures on a chosen grid.

An externally clocked \emph{allocation flow} consists of $\nu_0=\delta_{q_0}$,
$\nu_t\in\mathcal P_{m_t}(K_t)$, and a probability $\alpha_t(\cdot\mid q)$ supported on $\mathcal A_t$ at each atom of $\nu_t$, such that
\begin{equation}\label{eq:flow}
 \lambda_{t+1}=\int\sum_a\alpha_t(a\mid q)\Lambda_t(q,a)\nu_t(dq),
 \qquad \nu_{t+1}\cx\lambda_{t+1}\quad(0\le t<n).
\end{equation}
The measure $\lambda_{t+1}$ is the actual acquisition mixture proposed by the retained states and actions. It is not the law of a separately run exact filter. Define
\begin{equation}\label{eq:phi}
 \Phi_{\ext}(\mathbf m)=\inf_{\text{flows }\eqref{eq:flow}}
 \left\{\int g\,d\nu_n-B_n^*\right\}.
\end{equation}
All infima below allow arbitrarily close competitors; no unproved attainment assertion is implicit.

For a proposed autonomous label set $Z$, initial label $z_0$, and stopping subset $D$, put
$p_{t,z}\ge0$, $\sum_zp_{t,z}=1$, $q_{t,z}\in K_t$, and fix one action probability $\alpha_z(a)$ for every $z\notin D$. Require
\begin{equation}\label{eq:occupancy}
 p_{0,z}=\mathbf1_{z=z_0},\quad q_{0,z_0}=q_0,\quad
 p_{t,z}=0\ (t<n,z\in D),\quad p_{n,z}=0\ (z\notin D).
\end{equation}
In addition require $p_{t,z}\alpha_z(a)=0$ whenever $a\notin\mathcal A_t$. All references to feasible occupancies below include this legality constraint. The same label may have different \emph{conditional predictive states} $q_{t,z}$ at different cuts; these are proposed law parameters, not time-varying registers readable by the machine. For each $z\notin D,a$, set
\begin{equation}\label{eq:domination}
 \eta_{z,a}=\sum_{t=0}^{n-1}p_{t,z}\alpha_z(a)J_t(q_{t,z},a),\qquad
 h_{t,z,a}=\frac{d[p_{t,z}\alpha_z(a)J_t(q_{t,z},a)]}{d\eta_{z,a}}.
\end{equation}
If $\eta_{z,a}=0$, its contribution is zero. Otherwise $\sum_t h_{t,z,a}=1$ almost everywhere. For arbitrary $\phi_{t,j}\in\aff(K_{t+1})$, define
\begin{equation}\label{eq:simdual}
\begin{split}
 \sum_{t=0}^{n-1}\sum_{j\in Z}p_{t+1,j}\phi_{t,j}(q_{t+1,j})
 \ \le\sum_{z\notin D}\sum_a\int
 \max_{j\in Z}\left\{
  \sum_{t=0}^{n-1}h_{t,z,a}(y)\phi_{t,j}(F_t(q_{t,z},a,y))
 \right\}\eta_{z,a}(dy).
\end{split}
\end{equation}
Call the proposed occupancies \emph{simultaneously allocatable} if \eqref{eq:simdual} holds for every such family. Constants among the $\phi_{t,j}$ enforce the masses as well as the barycenters. Finite linear combinations of the determining test coordinates suffice. Define
\begin{equation}\label{eq:phiaut}
 \Phi_{\aut}(n,M)=\inf_{\substack{|Z|\le M,\ \eqref{eq:occupancy}\\
                    \text{all inequalities }\eqref{eq:simdual}}}
       \left\{\sum_zp_{n,z}g(q_{n,z})-B_n^*\right\}.
\end{equation}
No observer transition kernel occurs among these variables.

\begin{theorem}[Causal allocation and exact resource--risk frontier]\label{thm:completion}
For the prepared instruments and terminal task of Section~\ref{sec:experiment}, the following assertions hold.

\textup{(a)} A finite-support externally clocked flow is realizable by a causal observer with the indicated phase alphabets if and only if it satisfies \eqref{eq:flow}. Equivalently, for every step and every collection $\phi_1,\ldots,\phi_m\in\aff(K_{t+1})$, its proposed atoms $(p_j,q_j)$ satisfy
\begin{equation}\label{eq:localdual}
 \sum_{j=1}^m p_j\phi_j(q_j)
 \le\int\max_{1\le j\le m}\phi_j(w)\,\lambda_{t+1}(dw).
\end{equation}
Every failed contraction has a strict witness of this finite polyhedral form.

\textup{(b)} Proposed labelled occupancies and action probabilities are realizable by one autonomous tuple of Definition~\ref{def:machines}, stopping exactly at $n$, if and only if \eqref{eq:occupancy} and all simultaneous inequalities \eqref{eq:simdual} hold. Whenever they fail, one finite family of affine tests strictly separates the proposed moments from all phase-independent report allocations.

\textup{(c)} With the same preparation, task and legal actions on both sides,
\begin{equation}\label{eq:exactrisk}
 R_{\ext}(\mathbf m)=B_n^*+\Phi_{\ext}(\mathbf m),\qquad
 R_{\aut}(n,M)=B_n^*+\Phi_{\aut}(n,M).
\end{equation}
Every feasible flow is synthesized with exactly its declared labels. The auxiliary predictive states used to specify the flow are programme data, not additional persistent state. The synthesized transition functions need only be measurable; finite-bit or efficient synthesis is not asserted.

\textup{(d)} Let $\nu_t$ be the retained-state laws of any calibrated observer, and $\lambda_{t+1}$ its actual pre-pooling law. If the terminal decision is optimal conditional on the retained label, then
\begin{equation}\label{eq:slack}
 R-B_n^*=\sum_{t=0}^{n-1}(A_t+J_t),
\end{equation}
where
\begin{align*}
 A_t&=\E\left[\sum_a\alpha_t(a\mid Z_t)
       \int V_{t+1}(w)\Lambda_t(Q_t,a)(dw)-V_t(Q_t)\right],\\
 J_t&=\int V_{t+1}\,d\nu_{t+1}-\int V_{t+1}\,d\lambda_{t+1}.
\end{align*}
Both terms are nonnegative. For an autonomous machine $\alpha_t(\cdot\mid z)=\alpha_z$ is the same kernel at every cut. A suboptimal terminal decision adds its nonnegative conditional decision slack.
\end{theorem}

The theorem does not assume a density, a stable update, finite-dimensional closure, or full task curvature. It is an exact geometric feasibility and optimization theorem on the declared compact predictive class, not a universal numerical formula for arbitrary Borel experiments. Its two dual tests differ: in \eqref{eq:simdual} the maximum is taken \emph{after summing over phases}. Moving that maximum inside the phase sum would allow a different programme at every unobserved phase and would be a strictly weaker test.

\begin{corollary}[Exact and approximate completion]\label{cor:zero}
A zero-excess observer exists with a specified interface and budget if and only if its corresponding feasible allocation flow has $A_t=J_t=0$ at every cut and zero terminal decision slack. For every $\epsilon>0$, an observer with excess at most $\epsilon$ exists if and only if a feasible flow has $\sum_t(A_t+J_t)\le\epsilon$. The infimal excess is zero if and only if such flows exist for every positive $\epsilon$.
\end{corollary}
\begin{proof}
Apply the realization parts of Theorem~\ref{thm:completion} and the sum of nonnegative terms in \eqref{eq:slack}. The last statement separates an infimum of zero from attained exact sufficiency.
\end{proof}
For example, a piecewise-affine Bayes risk can have zero Jensen loss on a nontrivial cell. Its acquired dimension alone cannot force a positive lower. At the other extreme, Section~\ref{sec:geometry} estimates the same Jensen defects by actual cell mass and distributional curvature.

\begin{proposition}[Checkpoint identity]\label{prop:checkpoint}
For a full-history policy $\pi$, let $\nu_n^\pi$ be its actual terminal predictive law and $B_n^\pi=\int g\,d\nu_n^\pi$. Then
\begin{equation}\label{eq:checkpoint}
 R_{\cp}(n,M)=\inf_\pi\left\{B_n^\pi+D_M^g(\nu_n^\pi)\right\},\qquad
 D_M^g(\lambda)=\inf_{\substack{\nu\in\mathcal P_M(K_n)\\\nu\cx\lambda}}
                      \left\{\int g\,d\nu-\int g\,d\lambda\right\}.
\end{equation}
In particular, the two policy-dependent terms are not minimized separately.
\end{proposition}
\begin{proof}
Conditional on a final label, affinity gives the conditional predictive barycenter and the best risk $g$ at that barycenter. Necessity and sufficiency of the terminal contraction follow from the one-step allocation lemma below. Full-history independent randomization may be conditioned out for this checkpoint relaxation; alternatively include its kernels in $\pi$ and the displayed infimum is unchanged. For fixed finitely many readouts, assigning each predictive state to a least conditional loss readout is Borel and no worse than a randomized assignment. Optimizing the readouts therefore also gives the equivalent partition definition of $D_M^g$. This purification concerns a terminal encoder, not a stationary finite controller.
\end{proof}

\section{Barycentric allocation and the deployed law}\label{sec:allocation}
We give the finite-target form of the classical convex-order argument explicitly, since it is also the building block for phase-compatible implementation.

\begin{lemma}[Finite barycentric allocation]\label{lem:allocation}
Let $K$ be compact metrizable and convex as above, $\lambda\in\mathcal P(K)$, and $\nu=\sum_{j=1}^m p_j\delta_{q_j}$. The following are equivalent:
\begin{enumerate}
\item there are Borel functions $r_j:K\to[0,1]$, $\sum_jr_j=1$ $\lambda$-almost everywhere, such that
\begin{equation}\label{eq:moments}
 \int r_j\,d\lambda=p_j,\qquad
 \int r_j(w)f(w)\lambda(dw)=p_jf(q_j)\quad(f\in\aff(K));
\end{equation}
\item all inequalities \eqref{eq:localdual} hold;
\item $\nu\cx\lambda$.
\end{enumerate}
When the inequalities fail, finitely many determining affine coordinates and $m$ affine expressions provide a strict witness. Atomic acquired laws are allowed, and randomization is not silently replaced by deterministic cells.
\end{lemma}
\begin{proof}
Consider the simplex of allocations $r=(r_1,\ldots,r_m)$ in $(L^\infty(\lambda))^m$. It is weak-star compact: it is contained in the unit ball, and positivity and the identity $\sum_jr_j=1$ are weak-star closed. Map it to all its masses and affine-coordinate moments. Each coordinate is an integral against an $L^1(\lambda)$ function, so the image is compact and convex in a product of real lines. If the target moments $(p_j,p_jq_j)$ are outside that image, separation in this locally convex product gives a continuous linear functional involving only finitely many coordinates. Combining the coordinates for each $j$ yields an affine $\phi_j$. The supremum of the separated functional over allocations is exactly
\[
 \int\max_j\phi_j(w)\lambda(dw):
\]
a measurable least maximizing index attains it. This proves equivalence of the first two assertions, including strict separation. A weak-star allocation has Borel versions on the standard Borel space; modify on one null set to keep it in the simplex everywhere. Matching the separating coordinates matches the barycenter, hence every continuous affine function.

Given \eqref{eq:moments}, generate $W\sim\lambda$ and choose label $j$ with probability $r_j(W)$. Then $\E[W\mid j]=q_j$, where conditional expectation is understood through affine coordinates. Jensen's inequality gives $\nu\cx\lambda$. Conversely, for any affine family, convex order gives
\[
 \sum_jp_j\phi_j(q_j)\le\int\max_i\phi_i\,d\nu
 \le\int\max_i\phi_i\,d\lambda,
\]
which is assertion two. This also proves sufficiency without invoking a separate infinite-dimensional transport theorem. The argument is the finite-target martingale/convex-order mechanism of Strassen\cite{strassen}, not a new static comparison theorem.
\end{proof}

\begin{lemma}[Calibration under all incoming reports]\label{lem:calibration}
Suppose label $i$ has probability $p_i$ and conditional predictive state $q_i$. Draw its action from $\alpha_i$, execute the raw report kernel, and let $W=F_t(q_i,a,Y)$. If a next-label rule satisfies $\E[W\mid Z'=j]=q'_j$, then $q'_j$ is the actual predictive state conditional on $Z'=j$. In particular this applies to the kernel $r_j(W)$ of Lemma~\ref{lem:allocation} under the full acquired mixture $\lambda=\sum_{i,a}p_i\alpha_i(a)\Lambda_t(q_i,a)$.
\end{lemma}
\begin{proof}
The incoming label represents a conditional mixture of physical histories, not a physical preparation performed anew. Equation~\eqref{eq:affine_instrument} and predictive sufficiency imply that the conditional expectation of each determining future test, given $(i,a,Y)$, is its affine evaluation at $W$. Integrating this identity over \emph{all} incoming labels and reports sent to $j$ gives its evaluation at $q'_j$. The tests determine the predictive quotient. This is precisely the conditional state under the deployed observer. There is no comparison with the full-history posterior on a different controlled path.
\end{proof}

\begin{proof}[Proof of Theorem~\ref{thm:completion}(a), (c) for the external interface, and (d)]
For necessity, at each persistent cut of any observer take its actual conditional predictive states. Their weighted law is $\nu_t$, with support size at most $|Z_t|$. If several labels have the same state, average their action distributions using their conditional label probabilities. This gives the same acquired mixture $\lambda_{t+1}$. The next retained state is a conditional barycenter of the pre-pooling state, so Jensen's inequality gives \eqref{eq:flow}. This step records only marginals needed for subsequent predictive evolution; it does not assert equality of the original labelled trajectories after duplicate states are merged.

For sufficiency, assume a flow. Use its atoms as phase labels and its action probabilities as the programmed action kernel. Lemma~\ref{lem:allocation} provides $r_{t+1,j}$ for $\lambda_{t+1}\to\nu_{t+1}$. Programme
\[
 k_t(j\mid i,a,y)=r_{t+1,j}(F_t(q_{t,i},a,y)).
\]
The conditional-state parameters are constants in this phase's programme, not a persistently updated real vector. Lemma~\ref{lem:calibration} proves the prescribed occupancies and states by induction from $q_0$. Thus each feasible flow yields a legal observer of risk $\int g\,d\nu_n$. Conversely, replacing any observer's final readout by the best conditional readout gives this risk. Infimizing proves the first equality in \eqref{eq:exactrisk}.

For the telescope, let $Q_t$ denote the actual retained conditional state. Before pooling at the next cut,
\[
 \E V_{t+1}(W_{t+1})-\E V_t(Q_t)=A_t\ge0
\]
by the Bellman minimum. Calibration and concavity give
\[
 \E V_{t+1}(Q_{t+1})-\E V_{t+1}(W_{t+1})=J_t\ge0.
\]
Adding and telescoping yields \eqref{eq:slack}. The argument uses the machine's actual action distribution and remains valid at nonsmooth action switches. It neither compares independently optimized shadow controllers nor assumes Bellman-optimal actions for a stationary machine that reuses a label.
\end{proof}

\begin{remark}[Why the dual is operational]\label{rem:dual}
An arbitrary terminal splitting of a preparation is not automatically obtainable from a fixed sensor. The right side of \eqref{eq:localdual} is evaluated under that sensor's actual $\lambda$, not under an unconstrained information-design experiment. Likewise, local contractions at different cuts do not authorize phase-dependent code. The simultaneous test below imposes precisely the shared-kernel restriction that ordinary one-step convex order omits.
\end{remark}

\subsection{Finite obstructions to a global risk target}
For the externally clocked problem, compactness upgrades the witness for one proposed flow to a finite obstruction valid for every flow at an unattainable risk level. This is not an interchange of an infimum with a pointwise dual supremum.

\begin{theorem}[Compact completion and finite obstructions]\label{thm:finite_obstruction}
Under the compactness and weak-continuity hypotheses of Theorem~\ref{thm:completion}, every finite external width profile has an optimal allocation flow and an optimal measurable observer. For each phase fix a countable family $h_{t,1},h_{t,2},\ldots$ dense in the continuous convex functions on $K_t$, each a finite maximum of finite affine combinations of the determining coordinates. Let $c_k$ be the minimum terminal risk subject to the atomic-width, action-marginal and acquisition constraints of \eqref{eq:flow}, but retaining only the first $k$ convex-order inequalities at each cut. Then
\begin{equation}\label{eq:finite_obstruction}
 c_k\uparrow R_{\ext}(\mathbf m).
\end{equation}
Consequently, if a number $r<R_{\ext}(\mathbf m)$ is proposed as a feasible risk, finitely many of these acquired-law inequalities, together with the resource and risk constraints, already exclude every proposed flow with risk at most $r$. This is a finite family of tests, not necessarily a single test that defeats every candidate. No effective bound on its size or computability is asserted for general compact Borel data.
\end{theorem}
\begin{proof}
Represent the action assignments by probabilities $\gamma_t$ on $K_t\times\mathcal A_t$ whose first marginal is $\nu_t$. The sets of probabilities with at most $m_t$ atoms are compact: a sequence of such probabilities has convergent subsequences of weights and atom locations after padding by zero weights. Marginal constraints are closed. The map
$\gamma_t\mapsto\lambda_{t+1}=\int\Lambda_t(q,a)\gamma_t(dq,da)$
is weakly continuous by weak continuity of each $\Lambda_t$ and finiteness of the action set. Thus the candidate space before imposing convex order is compact. Every convex-order inequality is closed, and terminal risk is continuous. Feasibility is nonempty: choose one legal action at every cut and replace each acquired law by its single barycenter. Hence a risk minimizer exists and is synthesized by Theorem~\ref{thm:completion}.

For completeness, the required countable polyhedral family exists. The closed convex epigraph of a continuous convex function on compact $K_t$ can be separated from each point lying strictly below it by a continuous affine functional of the ambient countable product and the height coordinate. The height coefficient cannot vanish, because the epigraph contains a point with the same base coordinate. The resulting affine minorant uses only finitely many coordinates. A finite subcover of $K_t$ gives a finite maximum approximating the function uniformly from below. Approximating its finitely many coefficients by rationals gives a countable uniformly dense collection, with constants included. Inequalities for this collection imply all continuous convex inequalities by uniform approximation.

The finite-constraint feasible sets are nonempty, compact and nested, so their minima $c_k$ are nondecreasing and at most the full optimum. Choose minimizers and a convergent subsequence in the compact candidate space. For every fixed test index, all sufficiently late members obey its inequality, and the limit does too. Thus the limit satisfies all convex-order constraints, while its terminal risk is $\lim_kc_k$. This proves \eqref{eq:finite_obstruction}. If $r$ is strictly below that limit, some finite $k$ has $c_k>r$, yielding the asserted obstruction. Its proof does not bound $k$ and does not establish a finite computational representation of the carrier or kernels.
\end{proof}

\section{One programme across all phases}\label{sec:stationary}
\begin{lemma}[Simultaneous report allocation]\label{lem:simultaneous}
For fixed $(p_{t,z},q_{t,z},\alpha_z)$ satisfying \eqref{eq:occupancy}, the inequalities \eqref{eq:simdual} are equivalent to the existence of Borel probabilities $k(j\mid z,a,y)$, independent of $t$, such that, for every $t<n$, label $j$ and $f\in\aff(K_{t+1})$,
\begin{equation}\label{eq:sim_moments}
 p_{t+1,j}f(q_{t+1,j})=
 \sum_{z\notin D,a}p_{t,z}\alpha_z(a)
 \int k(j\mid z,a,y)f(F_t(q_{t,z},a,y))J_t(q_{t,z},a)(dy).
\end{equation}
The case $f=1$ is included.
\end{lemma}
\begin{proof}
For each $(z,a)$ use the finite dominating measure $\eta_{z,a}$ in \eqref{eq:domination}. In the finite product of spaces $(L^\infty(\eta_{z,a}))^{|Z|}$, take the weak-star compact convex set of simplex-valued functions. The map to the right sides of \eqref{eq:sim_moments} is continuous into the product of affine-coordinate moments: its test functions are
$h_{t,z,a}(y)f(F_t(q_{t,z},a,y))$, which are in $L^1(\eta_{z,a})$. If the target moments are not in this compact convex image, finite-coordinate separation gives a family $\phi_{t,j}$. For each incoming pair $(z,a)$ the support function of its allocation simplex is
\[
 \int\max_j\sum_t h_{t,z,a}(y)\phi_{t,j}(F_t(q_{t,z},a,y))\,\eta_{z,a}(dy).
\]
Summing over $(z,a)$ is exactly the right side of \eqref{eq:simdual}. A Borel least maximizing index realizes this support function. Separation proves both directions and gives a strict finite witness of infeasibility. Borel versions and zero-measure pairs are treated as in Lemma~\ref{lem:allocation}.
\end{proof}

\begin{proof}[Proof of Theorem~\ref{thm:completion}(b) and the autonomous part of (c)]
An actual autonomous tuple supplies its cut occupancies and conditional states. Its single transition kernel satisfies \eqref{eq:sim_moments} by conditioning on the incoming label, action and report. The stop rule gives \eqref{eq:occupancy}. Thus the simultaneous inequalities are necessary.

Conversely, Lemma~\ref{lem:simultaneous} constructs one kernel $k(j\mid z,a,y)$ for all phases. Use it with the proposed, phase-independent $\alpha_z$. Starting from $z_0$, induction in \eqref{eq:sim_moments} gives exactly $p_{t,z}$ and $q_{t,z}$, by the all-incoming calibration argument. The induction does not require the machine to know $q_{t,z}$ or $t$: those objects appear only in its offline feasibility certificate, while the executed kernel is the single $k$. The zero occupancies in \eqref{eq:occupancy} ensure that it never reaches a stopping label before $n$ and reaches one at $n$ almost surely. Choose a best readout at each positive-mass terminal label. This gives the stated risk and uses only $Z$. Taking infima proves the second equality in \eqref{eq:exactrisk}.
\end{proof}

\begin{remark}[A strict distinction between the tests]\label{rem:gluing}
Suppose incoming label $1$ and a fixed action see the same deterministic report in two phases, and both next predictive carriers consist of a singleton. A proposal sends that report to label $1$ with probability one in the first phase and to label $2$ with probability one in the second. Each separate barycentric allocation is feasible. Set $\phi_{0,1}=\phi_{1,2}=1$ and the other two constants to zero, with equal incoming phase masses. The left side of the simultaneous inequality is $2$, and its right side is $1$. Thus no single kernel implements the proposal. The obstruction is not a failed geometric covering; it is incompatible reuse of one programme location.
\end{remark}

\begin{proposition}[Exact conversion of timing conventions]\label{prop:clock}
For every experiment in the main theorem,
\begin{equation}\label{eq:conversion}
 R_{\ext}(n,M)\le R_{\aut}(n,M),\qquad
 R_{\aut}(n,1+\textstyle\sum_{t=1}^n m_t)\le R_{\ext}(\mathbf m).
\end{equation}
For the second inequality the source observer's phase label is included in the autonomous label. Suppose additionally that, for each action $a$, every $J_t(q,a)$ is mutually absolutely continuous with one common probability $\zeta_a$ on reports, uniformly in the sense of null sets over $t,q$. Then every internally exact-$n$ observer has disjoint positive-probability label supports at all $n+1$ cuts, and
\begin{equation}\label{eq:clockidentity}
 R_{\aut}(n,M)=
 \inf_{\substack{m_t\ge1\\1+\sum_{t=1}^n m_t\le M}} R_{\ext}(\mathbf m).
\end{equation}
Its class is nonempty exactly when $M\ge n+1$. No uniform likelihood-ratio lower is required.
\end{proposition}
\begin{proof}
An external observer may ignore its clock and run an autonomous one, proving the first inequality. Conversely, attach $t$ to each phase label of an external observer; its initial label is separate and its final labels stop. This gives a stationary tuple on the disjoint union and the second inequality.

For disjointness, suppose a positive-mass label $z$ occurs at cuts $s<t$. Starting from that label, run the same stationary action and transition kernels against reference reports $\zeta_a$ and the same fresh randomization. For every finite continuation length, successive conditioning shows equivalence between actual and reference trace laws: the next-report density is positive and finite almost everywhere, and randomization is unchanged. Integrating over any incoming conditional physical state or history preserves the same null sets. From cut $t$, stopping after $n-t$ further calls has probability one, hence reference probability one. Starting from the reused label at cut $s$, the reference machine is identical, so the same event has actual probability one. It stops before $n$, a contradiction. The case $t=n$ follows directly from the stopping set. Thus the cut supports are disjoint. Their sizes obey $1+\sum_tm_t\le M$, giving the reverse inequality in \eqref{eq:clockidentity}. A chain with $n$ legal calls and one terminal label proves nonemptiness at $n+1$, irrespective of its quality of prediction.
\end{proof}
Report time stamps or a freely supplied stage index invalidate the disjointness argument, not the simultaneous allocation theorem. An autonomous machine with such reports is still governed by \eqref{eq:simdual}; it may be able to reuse labels. The conversion result therefore specifies which timing cost is a resource convention and which compatibility constraint survives that convention.

\section{Quantitative acquired geometry}\label{sec:geometry}
The exact criterion separates implementability from estimates on the objective. We now re-prove a regular quantitative consequence, retaining the all-incoming and clock hypotheses explicitly. This section's curvature and integer-allocation mechanism is inherited from the preceding article\cite{prior}; the exact criterion above is not inferred from it.

All carriers in this section are one compact convex $K\subset\R^d$ with nonempty affine interior. Fix constants comparing the chosen norm $\|\cdot\|$ to the Euclidean norm $|\cdot|$. Suppose every $\Lambda_t(q,a)$ has density at most $\rho$ relative to the displayed $d$-dimensional affine volume, for \emph{all} $q,a,t$. Suppose $g$ is Lipschitz and $\kappa$-strongly concave in Euclidean coordinates, $\kappa>0$. Assume $J_t$ is uniformly TV-continuous and
\begin{equation}\label{eq:regular_stability}
 W_1(\Lambda_t(q,a),\Lambda_t(q',a))\le\beta_t\|q-q'\|,
 \quad 0<\beta_t\le B,\qquad
 \prod_{j=t}^{t+b-1}\beta_j\le\theta<1
\end{equation}
for every complete block of a fixed length $b$. The same experimental constants are used for every horizon. Zero coefficients can be enlarged to positive sufficient coefficients; no claim of sharpness for $\beta_t$ is made.

\begin{theorem}[Regular acquired-geometry transfer]\label{thm:regular}
Under these hypotheses there are constants $0<c\le C<\infty$, depending only on $K,d$, the norm comparison, $\rho,\kappa,\Lip(g),B,b,\theta$, such that, for every $n\ge1$ and $M\ge1$,
\begin{equation}\label{eq:regular_ext}
 cM^{-2/d}\le R_{\cp}(n,M)-B_n^*
 \le R_{\ext}(n,M)-B_n^*\le CM^{-2/d}.
\end{equation}
If report laws also have the common null sets of Proposition~\ref{prop:clock}, then the autonomous class is nonempty exactly when $M\ge n+1$ and
\begin{equation}\label{eq:regular_aut}
 c(M-n)^{-2/d}\le R_{\aut}(n,M)-B_n^*
 \le C(M-n)^{-2/d}.
\end{equation}
All legal adaptive acquisition policies are included. Without the block-product assumption the same construction gives, with $M\ge n+1$ in the autonomous expression,
\begin{align}\label{eq:weighted_upper}
 R_{\ext}(n,M)-B_n^*&\le C_0M^{-p}\sum_{t=1}^n a_t,\\
 R_{\aut}(n,M)-B_n^*&\le C_0(M-n)^{-p}
         \left(\sum_{t=1}^n a_t^{1/(1+p)}\right)^{1+p},\notag\\
 p&=2/d,\qquad a_t=\Lip(g)\prod_{j=t}^{n-1}\beta_j.\notag
\end{align}
The autonomous upper does not need clock-neutral reports; that hypothesis is used only for its matching lower.
\end{theorem}
For a finite measure $\lambda$ and a partition $\mathcal P$, set
\[
 \mathcal J_v(\lambda,\mathcal P)=\sum_{C\in\mathcal P}
 \lambda(C)v(\bar\lambda_C)-\int v\,d\lambda.
\]
Zero-mass cells are omitted. This is the Jensen objective of the corresponding barycentric contraction. The following estimate includes singular acquired measures, although a clean power law needs an additional mass estimate.

\begin{lemma}[Boundary-valid acquired-mass curvature estimate]\label{lem:curvature}
Let $K\subset\R^d$ be compact, convex, and have nonempty interior. If $v$ is concave and Euclidean $L$-Lipschitz on $K$, then $-v$ has a globally convex $L$-Lipschitz extension $f$. For every finite measure $\lambda$ on $K$ and grid side $h>0$,
\begin{equation}\label{eq:curvbound}
 \mathcal J_v(\lambda,\mathcal P_h)
 \le C_d h^{2-d}\sum_C\lambda(C)
 \kappa_f(\overline B_{3\sqrt d h}(z_C)).
\end{equation}
If $\lambda$ is a probability law with density at most $\rho$, then for every integer $L_0\ge1$ there is a Borel partition with at most $L_0$ cells such that
\begin{equation}\label{eq:gridrate}
 \mathcal J_v(\lambda,\mathcal P)\le C_{K,d}\rho L L_0^{-2/d}.
\end{equation}
Neither an interior support margin nor differentiability of $v$ is required.
\end{lemma}
\begin{proof}
Put $f_0=-v$. At each $x\in\operatorname{int}K$ select all supporting slopes of $f_0$. Their Euclidean norms are at most $L$: directional difference quotients can be taken in every direction at an interior point. Take the supremum, on all $\R^d$, of the resulting supporting affine functions. This supremum is finite, convex, $L$-Lipschitz and agrees with $f_0$ on $\operatorname{int}K$; continuity gives equality on $K$. In particular, no assertion about arbitrary Lipschitz extensions preserving convexity is being used.

First suppose the extension is smooth. For a cell center $z$, subtract its tangent plane and call the resulting convex function $w$. Thus $w\ge0$ and $w(z)=0$. Put $R=2\sqrt d h$. For any $x$ in the cell, $B_{R/2}(x)\subset B_R(z)$. Writing $\operatorname{av}$ for the normalized spherical average, the submean inequality for $w$, radial monotonicity, and the divergence theorem give
\[
 w(x)\le \frac{C_d}{R^d}\int_{B_R(z)}w
 \le C_d\operatorname{av}_{\partial B_R(z)}w
 \le C_d R\operatorname{av}_{\partial B_R(z)}\partial_r w
 =C_dR^{2-d}\int_{B_R(z)}\Delta f.
\]
Here radial monotonicity gives $\int_{B_R}w\le (R/d)\int_{\partial B_R}w$.
The radial inequality is $w(z+R\omega)\le R\partial_r w(z+R\omega)$. This argument also applies in dimension one, with the boundary integral interpreted as the sum over two endpoints. If $b_C$ is the conditional barycenter, the affine tangent terms cancel and
\[
 \int_Cf\,d\lambda-\lambda(C)f(b_C)
 =\int_Cw\,d\lambda-\lambda(C)w(b_C)
 \le\lambda(C)\sup_Cw.
\]
Mollify $f$ on $\R^d$. The convex mollifications converge uniformly on bounded sets, and their positive Laplacian measures converge weakly to $\kappa_f$. The limsup of the mass of the radius-$R$ ball is bounded by the mass of the closed ball of radius $3\sqrt d h$. Summing the finitely many relevant cells proves \eqref{eq:curvbound}.

A Lipschitz convex function has uniformly bounded Laplacian mass on any fixed bounded region: choose a smooth nonnegative cutoff $\chi$ that equals one there and use
\[
 \int\chi\,d\kappa_f=-\int\nabla\chi\cdot\nabla f
 \le L\int|\nabla\chi|.
\]
For $h$ bounded above, the enlarged cell balls have bounded overlap, and $\lambda(C)\le\rho h^d$. Equation~\eqref{eq:curvbound} therefore gives $C_{K,d}\rho Lh^2$. Enclose $K$ in a fixed cube and subdivide each coordinate into $\lfloor L_0^{1/d}\rfloor$ intervals. There are at most $L_0$ cells. Since $\lfloor L_0^{1/d}\rfloor\ge L_0^{1/d}/2$, one has $h\le C_K L_0^{-1/d}$. This proves \eqref{eq:gridrate}, including $L_0=1$, by changing the fixed constant. The extension is used only in this proof, not as an observer state.
\end{proof}

\begin{lemma}[Integer allocation with a baseline label]\label{lem:integer}
Let $a_t\ge0$, $p>0$, and $m\ge1$. There are integers $L_t\ge1$ with $\sum_t L_t\le n+m-1$ and
\begin{equation}\label{eq:integer}
 \sum_t a_tL_t^{-p}\le 2^p m^{-p}
       \left(\sum_t a_t^{1/(1+p)}\right)^{1+p}.
\end{equation}
When all $a_t=0$ choose $L_t=1$. When $m=1$ the same choice is valid.
\end{lemma}
\begin{proof}
Otherwise put $w_t=a_t^{1/(1+p)}$, $A=\sum_tw_t$ and
$L_t=1+\lfloor(m-1)w_t/A\rfloor$. The total count is as stated. For $m=1$, $\sum_t a_t\le A^{1+p}$. For $m\ge2$ and $w_t>0$, $L_t\ge(m-1)w_t/A$ gives $\sum_ta_tL_t^{-p}\le A^{1+p}(m-1)^{-p}\le2^p A^{1+p}m^{-p}$. Zero-weight phases still have one real label.
\end{proof}

\begin{proof}[Proof of Theorem~\ref{thm:regular}]
At each phase the action set is common to all incoming states. The Wasserstein bound therefore implies
$\Lip(V_t)\le\beta_t\Lip(V_{t+1})$, by comparing the same action at two states before taking the finite minimum. Thus $\Lip(V_t)\le a_t$. Concavity was established independently from fixed-programme affine risks.

Starting from the preparation, choose a Bellman-minimizing action at every current label, partition the \emph{full incoming mixture} at the next cut, and use its cell barycenters. Lemma~\ref{lem:calibration} makes this a calibrated machine under its own law. Every pre-pooling law obeys the density bound, since the bound holds at all incoming mixture states and actions. Lemma~\ref{lem:curvature} with $L_t$ cells and the exact identity \eqref{eq:slack} give
$R-B_n^*\le C_0\sum_ta_t L_t^{-p}$. Choosing $L_t=M$ proves the external weighted upper. Choosing the integers of Lemma~\ref{lem:integer} with $m=M-n$ and using disjoint phase labels proves the autonomous weighted upper with total count $1+\sum_tL_t\le M$. No exact filter is executed outside these labels.

For $k=n-t$, block decomposition yields the explicit estimate
\[
 a_t\le\Lip(g)\max(1,B^{b-1})\theta^{\lfloor k/b\rfloor}.
\]
Hence both sums in \eqref{eq:weighted_upper} are uniformly bounded. This proves the two horizon-uniform uppers, including $m=1$.

For the lower, condition on the full history immediately before the final action/report. Every conditional final predictive law has density at most $\rho$, so its mixture under any randomized adaptive policy does too. If $S$ is this full-history state and a final label $Z$ takes at most $L$ values, its best forecast has conditional state $Q=\E[S\mid Z]$. Strong concavity implies
\[
 \E g(Q)-\E g(S)\ge\frac{\kappa}{2}\E|S-Q|^2.
\]
The $L$ balls of radius $r=(2\rho v_dL)^{-1/d}$ about the possible values of $Q$ cover mass at most $1/2$. Thus $\E|S-Q|^2\ge r^2/2$. This also covers a randomized label assignment: $Q$ still has at most $L$ values. Since the full-history risk of this policy is at least $B_n^*$, the same lower holds above $B_n^*$. Use $L=M$ for checkpoint and external observers. For an autonomous observer, Proposition~\ref{prop:clock} gives $L\le M-n$. This completes the proof. The lower concerns acquired mass under every admissible policy, not a packing of a reachable set.
\end{proof}

\begin{proposition}[Finite strata with unnormalized mass]\label{prop:strata}
Replace the density hypothesis by a finite report-readable decomposition of each $\Lambda_t(q,a)$ into submeasures supported on fixed compact convex affine sets $K_i\subset K$, of dimensions $d_i\le r$, $r\ge1$. Positive-dimensional submeasures have densities at most $\rho_i$ in their own affine volumes; zero-dimensional types are atoms. Suppose every full-history terminal law has an $r$-dimensional submeasure of actual mass at least $w>0$ and density at most $\rho_*$. Retain the task and block hypotheses on $K$. Then \eqref{eq:regular_ext}--\eqref{eq:regular_aut} hold with exponent $2/r$. The lower constant retains a factor
\begin{equation}\label{eq:stratum_mass}
 \kappa\,w^{1+2/r}\rho_*^{-2/r}.
\end{equation}
The autonomous lower again requires common report null sets.
\end{proposition}
\begin{proof}
For a phase allocation $L\ge2s$, where $s$ is the number of types, assign one cell to every atomic type and at least $\lfloor L/s\rfloor$ cells to each positive-dimensional type, reducing counts if necessary so the total is at most $L$. Restrict the curvature lemma to each $K_i$ and each unnormalized component. Since $d_i\le r$, its loss is at most a fixed multiple of $\Lip(V_t)L^{-2/r}$. Type is available from the current report, so the pair (type, cell) is an executable label rule. All incoming edges of that pair are pooled together. The resulting barycenter is in $K_i$. For $L<2s$, one global cell costs at most $\Lip(V_t)\operatorname{diam}(K)$; enlarge the constant by $(2s)^{2/r}$. Integer allocation and block summation proceed as before.

For the lower choose radius $r_0=(w/(2\rho_*v_r L))^{1/r}$. The union of $L$ ambient balls intersects the affine component in volume at most $Lv_r r_0^r$, and therefore covers at most $w/2$ of that submeasure. The remaining mass contributes at least $(w/2)r_0^2$ to squared distortion. Strong concavity gives \eqref{eq:stratum_mass}. No conditional normalization removes $w$, and no mass from other strata is subtracted.
\end{proof}

\section{Two raw-kernel realizations and conventional tasks}\label{sec:realizations}
These verifications are consequences of the general theorems, never their premises. Their kernel mechanisms were established in the preceding article\cite{prior}; we include complete proofs and specify ordinary prediction tasks.

\begin{lemma}[Positive-instrument coupling]\label{lem:coupling}
Let states be probability vectors in $\ell^1$ or finite-dimensional density matrices in trace norm. Let $C$ be a positive trace-preserving affine preparation channel with contraction coefficient $\gamma$. Let $A_y$ be a positive linear instrument density with $\int A_y\,\zeta(dy)$ trace preserving. For $q=C(s)$ put $u_y=A_y(q)$, $r_y=\tr u_y$, and $F_y=u_y/r_y$ on positive-mass reports. Then
\begin{equation}\label{eq:coupling}
 W_1(\law F_Y,\law F'_{Y'})\le3\gamma\|s-s'\|,
 \qquad\TV(\law Y,\law Y')\le\frac{\gamma}{2}\|s-s'\|.
\end{equation}
\end{lemma}
\begin{proof}
The positive and negative parts of a signed vector or Hermitian matrix $v$ give
\[
 \int\|A_y(v)\|\,\zeta(dy)\le\|v\|,\qquad
 \int|\tr A_y(v)|\,\zeta(dy)\le\|v\|.
\]
Couple common report mass $\min(r_y,r'_y)$ at the same $y$. There
\[
 \min(r_y,r'_y)\|F_y-F'_y\|\le\|u_y-u'_y\|+|r_y-r'_y|.
\]
The state space has diameter at most two, so the remaining report mass costs at most $\int|r_y-r'_y|\,d\zeta$. Summing proves the Wasserstein bound; the second integral directly gives the TV bound. These are sufficient coefficients, not claimed optimal ones.
\end{proof}

\subsection{Gaussian hidden-state sensing}
There are $d+1$ hidden states, a fixed finite action family, and known stochastic matrices $T_t^a$ with every entry at least $\tau>0$. A call applies $T_t^a$ and then reports
\[
 Y\mid X'=j,a\sim N(\mu_j^a,\Sigma_a),\qquad Y\in\R^{m_{\rm obs}},\quad m_{\rm obs}\ge d.
\]
The parameter set is compact: the means are uniformly bounded, the eigenvalues of $\Sigma_a$ are uniformly bounded above and below, and the least eigenvalue of
$B_a\Sigma_aB_a^T$ is uniformly positive, where $B_a$ has rows $(\mu_j^a-\mu_0^a)^T\Sigma_a^{-1}$. At each phase at least one legal $T_t^a$ is invertible. Put
$\gamma_t=\max_a\frac12\max_{i,j}\sum_k|T_t^a(i,k)-T_t^a(j,k)|$.

\begin{theorem}[Categorical filtering]\label{thm:gaussian}
If $\beta_t=3\gamma_t$ obeys \eqref{eq:regular_stability}, all conclusions of Theorem~\ref{thm:regular} hold for the conventional loss
\[
 L(X_n,u)=\|e_{X_n}-u\|_2^2,\qquad u\in\Delta_d,
\]
with predictive carrier $\Delta_d$. The full-history baseline is the usual optimal categorical filtering Bayes risk in this controlled experiment. Constants are uniform in $n,M$ on the stated parameter class. Every Gaussian report, including its tails, is a paid acquisition.
\end{theorem}
\begin{proof}
For incoming $p$ set $q=pT_t^a$; then $\tau\le q_j\le1-d\tau$. With the Gaussian emission densities $f_j$, the raw joint kernel gives
\[
 r(q,y)=\sum_jq_j f_j(y),\qquad u_j=F_j(p,a,y)=q_j f_j(y)/r(q,y).
\]
Unnormalized finite continuation probabilities are linear in $p$, proving the instrument mixture identities on the whole simplex. All report densities are positive on $\R^{m_{\rm obs}}$, and therefore equivalent to one fixed nonsingular Gaussian probability reference for each action, at every phase.

The actual posterior law has uniformly bounded density even at the simplex boundary. Indeed its logits satisfy
\[
 z_j=\log(u_j/u_0)=(B_ay)_j+(c_a)_j+\log(q_j/q_0),\qquad
 (c_a)_j=-\tfrac12(\mu_j^T\Sigma_a^{-1}\mu_j-\mu_0^T\Sigma_a^{-1}\mu_0).
\]
Their law is a finite mixture of nonsingular $d$-dimensional Gaussians, with uniform bounds on means and covariance eigenvalues. Hence its density obeys
$h_q(z)\le C\exp(-c|z|^2+C|z|)$. The softmax map has Jacobian $\prod_{j=0}^du_j$, and
\[
 (\textstyle\prod_j u_j)^{-1}
 =(1+\sum_{j=1}^d e^{z_j})^{d+1}e^{-\sum_j z_j}
 \le C_d e^{C_d|z|}.
\]
Consequently $h_q(\operatorname{logit}u)/\prod_j u_j$ is bounded uniformly on the entire simplex interior. The boundary has zero acquired mass; no clipping or conditioning was used.

The row coefficient gives $\|pT-p'T\|_1\le\gamma_t\|p-p'\|_1$: decompose the difference into equal positive and negative masses and bound the distance of two row mixtures. Lemma~\ref{lem:coupling} proves report continuity and the required Wasserstein estimate. This also gives weak continuity of the acquired kernel.

For predictive separation, distinct Gaussian means give linearly independent emission densities: divide a purported relation by the common quadratic Gaussian factor and restrict the resulting exponentials to a line with distinct slopes. Successive limits force all coefficients to vanish. The next-report law thus identifies $q$, and an invertible legal $T_t^a$ identifies $p$. Conversely $p$ determines all future tests, so this is the predictive quotient rather than a redundant parameterization.

Finally $\ell(p,u)=1-2p\cdot u+\|u\|_2^2$ is affine in $p$ and has Bayes risk $g(p)=1-\|p\|_2^2$. It is Lipschitz and strongly concave on the affine simplex. Apply Theorem~\ref{thm:regular}. The proof optimizes over the full declared adaptive policy class; it does not infer a strict feedback advantage merely from the existence of several sensors.
\end{proof}
For instance, two-state invertible transition matrices with row coefficients repeating $0.8,0.8,0.02$ have individual certified coefficients as large as $2.4$, but every three-step product is $0.3456<1$. Weak mixing by itself is not being substituted for this explicit block condition.

\subsection{Noncommuting finite-dimensional instruments}
Let $D\ge2$, $d=D^2-1$, and take an orthonormal real basis $H_1,\ldots,H_d$ of traceless Hermitian matrices. Choose $s>0$ with $s\sum_i\|H_i\|_{\rm op}<1$. With uniform probability measure on the report cube $[-1,1]^d$, put
\[
 E_y=I+s\sum_i y_iH_i,\qquad
 \mathcal I_{t,y}^a(\rho)=E_y^{1/2}\mathcal C_t^a(\rho)E_y^{1/2},
\]
where
\[
 \mathcal C_t^a(\rho)=(1-\gamma_t^a)\tau_t^a+
                    \gamma_t^a U_a\rho U_a^*,\qquad
 0<\gamma_-\le\gamma_t^a\le\gamma_+<1,\quad \tau_t^a\succeq\tau I.
\]
The affine channel extends linearly by multiplying its first term by $\tr\rho$. It is a paid mixture of preparation and a unitary channel. Since $\int E_y\,2^{-d}dy=I$, the instrument is positive and normalized. No common eigenbasis of the effects is assumed.

\begin{theorem}[Prediction of an informationally complete measurement]\label{thm:quantum}
If $\beta_t=3\max_a\gamma_t^a$ satisfies the block condition, Theorem~\ref{thm:regular} holds on the density-matrix carrier for Brier prediction of the outcome of a fixed finite informationally complete POVM. Affine volume means Lebesgue volume in the displayed orthonormal traceless chart, and the Wasserstein metric uses trace norm. The constants need not be uniform in $D$. The terminal measurement is paid when included in the physical acquisition count.
\end{theorem}
\begin{proof}
Let $q=\mathcal C_t^a(\rho)$, $r_y=\tr(E_yq)$ and
$\sigma_y=E_y^{1/2}qE_y^{1/2}/r_y$. Uniform positive lower bounds on $q,E_y$ give equivalent report supports on the whole cube. Linearity of the unnormalized quantum instruments proves mixture closure. The coupling lemma supplies the Wasserstein coefficient $3\gamma_t^a$ and TV continuity.

For the density calculation, at positive definite $q,\sigma$ define
\[
 A=q^{-1/2}(q^{1/2}\sigma q^{1/2})^{1/2}q^{-1/2}.
\]
This is the unique positive definite solution of $AqA=\sigma$. The inverse of the normalized congruence map on the trace-$D$ positive hyperplane is
\begin{equation}\label{eq:quantum_inverse}
 E=\frac{DA^2}{\tr A^2}.
\end{equation}
Since $A$ is positive definite, $E^{1/2}=\sqrt{D/\tr A^2}\,A$, proving the substitution in both directions. The map and its inverse are smooth; their differentials on the trace hyperplanes are inverse. The affine map $y\mapsto E_y$ has full rank. The permitted $q$ and closed cube lie in a compact subset of the positive domains, so the absolute Jacobian of $y\mapsto\sigma_y$ has a positive uniform lower bound. The bounded report density and one-to-one change of variables prove a uniform acquired-density upper. Cube boundaries are null, and no multiplicity factor is hidden.

The report density $1+s\sum_i y_i\tr(H_iq)$ identifies all traceless coordinates of $q$; the affine channel is injective since $\gamma_t^a>0$. Thus the density matrix is exactly the predictive quotient for the declared tests. It is derived from the raw instruments, not supplied to the finite observer as a register.

For a concrete fixed POVM take
$G_{i,\pm}=(I\pm eH_i)/(2d)$ with
$0<e<1/\max_i\|H_i\|_{\rm op}$. These effects sum to $I$. Its outcome probabilities $P_{i,\pm}(\rho)=\tr(G_{i,\pm}\rho)$ form an injective affine map. Under ordinary Brier loss the Bayes risk is $1-\|P(\rho)\|_2^2$. Its strong-concavity constant is twice the squared least singular value of the linear part of $P$ on the traceless chart. This is a probability-forecast task for an actual fixed measurement, not an assertion that every discrimination loss has full-dimensional curvature. Theorem~\ref{thm:regular} now applies.
\end{proof}
As a finite-stratum extension, replace a fixed fraction $\vartheta<1$ of calls by a flagged preparation of one of finitely many pure states and use the displayed continuous instrument otherwise. A subsequent positive channel sends pure states back into the interior. The type is visible in the current report; the continuous terminal submeasure has its actual mass $1-\vartheta$. Provided the corresponding coupling block bound holds, Proposition~\ref{prop:strata} applies with its unnormalized mass factor. This verifies genuine rank changes without asserting that arbitrary or unobserved quantum strata satisfy that proposition.

\section{Blind transport and singular acquisition}\label{sec:singular}
The exact criterion also applies when every displayed one-step stability coefficient is one, and when no affine density exists. This gives a sharp law different from the smoothing law of Section~\ref{sec:geometry}.

Let $X$ have a specified prepared law $\mu$ on a compact subset $E$ of a finite-dimensional Euclidean space. A protocol consists of one paid capture and $n-1$ paid blind calls. Only the action \emph{capture} is legal initially; it reports $X$ exactly. At subsequent nonterminal phases only \emph{advance} is legal and its report is a fixed symbol. The physical variable $X$ remains unchanged. After exactly $n$ calls the observer predicts $X$ under squared loss. The terminal scored outcome records $X$ after the decision, so bounded continuous functions of that outcome are executable terminal tests. It is unavailable before the prediction; an additional physical audit call, when needed, is counted separately. Thus the capture/advance interface is part of the experiment, not an uncounted clock supplied to the observer.

Define the actual-law quantization error
\[
 q_L(\mu)=\inf_{c_1,\ldots,c_L\in\operatorname{conv}(E)}
               \int\min_j\|x-c_j\|_2^2\mu(dx),\qquad L\ge1.
\]
The centers may be restricted to the compact convex hull, so a minimizer exists. The countable future tests on $E$ give the compact predictive carrier $\mathcal P(E)$; the capture law on this carrier is the pushforward of $\mu$ by $x\mapsto\delta_x$. Blind updates are the identity. Affine mixture closure holds for both raw kernels. No ambient smooth parameterization of $\mu$ is used.

\begin{theorem}[Exact blind-transport law]\label{thm:blind}
For this experiment, every external width profile $\mathbf m=(m_1,\ldots,m_n)$ satisfies
\begin{equation}\label{eq:blind_external}
 R_{\ext}(\mathbf m)=q_{\min_t m_t}(\mu).
\end{equation}
The stationary internally exact-$n$ class is nonempty precisely when $M\ge n+1$, and in that range
\begin{equation}\label{eq:blind_internal}
 R_{\aut}(n,M)=q_{\lfloor(M-1)/n\rfloor}(\mu).
\end{equation}
In particular, these conclusions hold for singular continuous and countably atomic acquired laws. They are assertions for the declared protocol and task, not for an experiment permitting repeated captures.
\end{theorem}
\begin{proof}
Fix a cut $t\ge1$. All future reports are constant, and fresh transition randomness is independent of $X$ conditional on the current label $Z_t$. Hence the terminal prediction $U$ satisfies the Markov relation $X-Z_t-U$. Conditional squared-loss decomposition gives
\[
 \E\|X-U\|^2\ge \E\|X-\E[X\mid Z_t]\|^2\ge q_{|\supp Z_t|}(\mu).
\]
The last inequality remains valid for randomized encoders: the conditional means give at most $|\supp Z_t|$ centers, and pointwise least squared distance is no larger than the encoder's conditional average. This lower holds at every cut. Conversely, quantize $X$ once using $L=\min_t m_t$ centers and carry the resulting label unchanged through all blind calls. The external clock supplies the legal phase and terminal request. This proves \eqref{eq:blind_external}.

For the autonomous machine, its initial label cannot occur later with positive probability. Its action must be capture initially and advance at every later nonterminal phase; the same stationary action kernel cannot do both. Nor can a nonterminal label be a stopping label. Between cuts $1$ and $n$, the reports are identical and the label dynamics are stationary. If one positive-probability label occurred at two different such cuts, its law of future stopping time would be the same from both, contradicting exact stopping at $n$. Thus these $n$ cut supports and the initial label are disjoint. Their sizes $L_t\ge1$ satisfy $1+\sum_{t=1}^nL_t\le M$, so one is at most $\lfloor(M-1)/n\rfloor$. The cut lower just proved gives \eqref{eq:blind_internal}. To attain it, use one initial capture label and $n$ disjoint copies of an optimal $L$-point quantizer label. The first copy is selected by the capture report, every blind transition advances its phase copy without changing its quantizer index, and the last copy stops and emits the center. The total number of states is $1+nL\le M$. The same construction with $L=1$ proves nonemptiness at $n+1$, and the disjoint-support argument proves necessity.
\end{proof}

This theorem is also an explicit solution of the allocation frontier in Theorem~\ref{thm:completion}. At capture, the law of conditional distributions is a barycentric contraction of the actual law of $\delta_X$. At every later cut a further contraction can only lose information. Carrying the chosen conditional distribution without further pooling makes all later Jensen defects zero, although the blind update has no strict contraction at all. Thus nonsummable bounds on generic Lipschitz sensitivities do not preclude exact causal completion. Conversely, the narrowest retained cut gives a lower unrelated to an artificial terminal density.

\begin{corollary}[Fractal and accumulating atomic laws]\label{cor:singular_rates}
\textup{(i)} Let $\mu$ be the equal-weight two-map Cantor law with contraction $r\in(0,1/2)$ and set $s=\log2/\log(1/r)$. Then
\[
 q_L(\mu)\asymp_r L^{-2/s}.
\]
\textup{(ii)} Let $\alpha>0$, $x_j=2^{-j}$, and
\[
 \mu\{x_j\}=(1-2^{-\alpha})2^{-\alpha(j-1)},\qquad j\ge1.
\]
On the compact support $\{0\}\cup\{x_j:j\ge1\}$ one has
\[
 q_L(\mu)\asymp_\alpha 2^{-(\alpha+2)L}.
\]
Substitution into \eqref{eq:blind_external}--\eqref{eq:blind_internal} gives matching causal laws, with the actual prepared masses unchanged.
\end{corollary}
\begin{proof}
For (i), a level-$k$ cylinder has diameter $r^k$ and mass $2^{-k}$. With $k=\lfloor\log_2L\rfloor$, its $2^k$ cylinder representatives give $q_L\le r^{2k}\le C_rL^{-2/s}$. The gaps separating neighboring cylinders at their first distinct level imply
$\mu(B(x,\rho))\le C_r\rho^s$ for every $x$ and $0<\rho\le1$: choose $k$ with $r^{k+1}<\rho\le r^k$; an interval of length $2\rho$ meets only a bounded number, depending on $r$, of level-$k$ cylinders. For $L$ centers choose $\rho=(2C_rL)^{-1/s}$, after increasing the constant for the finitely many coarse radii. At least half the acquired mass is outside their balls. This gives $q_L\ge\rho^2/2$.

For (ii), retain the first $L-1$ atoms exactly and place the remaining center at zero. Then
\[
 q_L\le\sum_{j\ge L}(1-2^{-\alpha})2^{-\alpha(j-1)}2^{-2j}
       \le C_\alpha 2^{-(\alpha+2)L}.
\]
The $L+1$ intervals centered at $x_j$ with radii $x_j/8$, $1\le j\le L+1$, are pairwise disjoint. At least one contains no one of $L$ proposed centers. Its atom alone contributes at least
$c_\alpha2^{-(\alpha+2)(L+1)}$ to the squared error. This lower is uniform over centers and proves the claim.
\end{proof}
The support in (ii) accumulates, and its effective finite-resolution curve is exponential rather than the power predicted by an ambient one-dimensional chart. Its proof uses the actual atom masses. The result does not assert a corresponding formula for arbitrary singular laws: for those, Theorem~\ref{thm:blind} retains the exact quantity $q_L(\mu)$.

\section{Causal simulation and finite execution}\label{sec:resources}
The preceding cardinalities count persistent labels. A programme describing real-valued functions or arbitrarily long tables is not free bit complexity. We separate statistical comparison, numerical execution, and one finite class where synthesis is effective.

\begin{proposition}[Typed common-task transfer]\label{prop:simulation}
Let $\mathfrak E$ and $\mathfrak F$ be prepared experiments with the same bounded scored task, $0\le L\le H$. Suppose a parameter-independent causal simulator implements every legal target strategy in a specified class, including its action interface, reports, stopping rule, and task readout. Assume it uses at most $R$ internal labels, transforms the target acquisition bound $N$ to a declared physical bound $C(N)$, and gives total variation at most $\varepsilon$ between the complete stopped and scored laws, uniformly in the compared preparations and strategies. Then
\begin{equation}\label{eq:simulation}
 \mathcal R_{\mathfrak E}(C(N),RM)
 \le \mathcal R_{\mathfrak F}(N,M)+H\varepsilon.
\end{equation}
Two such simulators compose with the product of their internal cardinalities, composition of their acquisition/physical-time maps, and the sum of their complete-law TV defects. A reverse risk lower requires a reverse certificate covering every strategy in the class to which the lower is applied.
\end{proposition}
\begin{proof}
For each target observer store the pair of its current label and the simulator label. The causal strategy lift supplies only legal actions, transforms reports at the promised completion times, and uses the same scored readout; the direct persistent cardinality is at most $RM$. A bounded loss changes by at most $H$ times total variation, so taking infima proves \eqref{eq:simulation}. For composition, insert the intermediate simulated law, use contraction of total variation under the next parameter-independent kernel, and then its own defect bound. The product-state and composed-call accounts follow from the direct implementation. Without a reverse strategy lift this construction says nothing about source strategies not obtained from a target observer, which proves the stated converse qualification.
\end{proof}
Programme descriptions concatenate, scratch workspaces add in a direct simultaneous implementation (or can be reused when their lifetimes are proved disjoint), and arithmetic and physical execution times add along the actual calls. An externally supplied stage index is not automatically supplied by a simulator of an internally timed task. The proposition compares total risks of a common task; subtracting two different full-history baselines would not preserve the inequality. Marginal report TV, an unscored transcript comparison, or a fixed-time simulator without stopping compatibility is insufficient. No large-deviation equivalence is inferred from this proposition.

\begin{proposition}[Certified execution error]\label{prop:numerical}
Consider an ideal finite observer for a task bounded by $H$ and at most $n$ physical transitions. Suppose an implementation has a coupling with the ideal process such that, while the complete pasts agree, the next physical/label transition fails to agree with conditional probability at most $e_t$. A bound under the ideal prefix law on the corresponding first-mismatch event is also sufficient. If the final readout perturbation changes the conditional task loss by at most $\omega(v)$, then
\begin{equation}\label{eq:numerical}
 R_{\mathrm{impl}}\le R_{\mathrm{ideal}}
    +H\min\{1,\sum_{t=0}^{n-1}e_t\}+\omega(v).
\end{equation}
An ideal observer using approximate Bellman selections adds its actual expected action slacks $\sum_t A_t$ as in \eqref{eq:slack}; these slacks need not be counted a second time in $e_t$.
\end{proposition}
\begin{proof}
Partition the event that the paths disagree by the first transition at which they do. Before that time both constructions have the ideal prefix, so the probability is at most $\sum_te_t$, capped at one. On the complementary event the scored outcome is the same apart from the asserted readout change. The loss bound proves the inequality. The action-slack statement is the exact telescope, not a stability estimate between separately optimized controllers.
\end{proof}
For the grids of Section~\ref{sec:geometry}, a coordinate error at most $b_t$ can change a cell only in slabs of total ideal acquired mass at most
$\min\{1,C_d\rho b_t/h_t\}$, for $b_t\le h_t$, plus a separately certified exceptional evaluation probability. Coordinates mean the fixed affine simplex or traceless-matrix chart used for volume and grid cells. These are unconditional ideal masses; conditioning on survival of earlier numerical comparisons would bias them. Certified kernel couplings, integration/barycenter routines, Bellman selection errors, report precision, programme bits, scratch space and time must be supplied to apply this observation. It is not an algorithm for arbitrary Borel inputs. Predictive-coordinate error used to choose a cell is distinct from a Brier probability-readout error, whose loss modulus on a bounded probability simplex is linear in its Euclidean size.

\begin{theorem}[Effective finite rational compilation]\label{thm:finite_compiler}
Suppose the physical state, action, report and terminal-decision sets are finite, with cardinalities $S,A,Y,U$, respectively. Preparation, phase-indexed raw transition probabilities and bounded losses are rational, and $n,M$ are fixed integers. For the stationary machine convention of Definition~\ref{def:machines}, including the phase-dependent physical legality constraints, nonemptiness and the optimum risk are decidable; whenever nonempty the optimum is attained by a machine with real algebraic row probabilities. No efficiency bound for the synthesis is asserted.

For every integer $Q\ge1$, this machine admits a rational, support-preserving rounding with denominator $Q$, at the same $M$, which is still legal and stops exactly at $n$, with risk increase at most
\begin{equation}\label{eq:compiler_error}
 H\min\{1,n(A+M)/Q\}.
\end{equation}
It can be stored using
\[
 O\big((MA+M^2AY)\log(Q+1)+M\log(U+1)+M\big)
\]
programme bits, in addition to the given model and a fixed interpreter. Its transition scratch space is
$O(\log(Q+1)+\log(M+1)+\log(A+1)+\log(Y+1))$ bits. Exact rational sampling has finite expected runtime; limiting each sample to $k$ rejection trials yields a bounded runtime and an additional risk error at most $2Hn2^{-k}$, without violating legality or the deadline. The program-size statement is a feasible upper, not an optimal description-length law.
\end{theorem}
\begin{proof}
Enumerate the finite choices of initial label, stopping subset and terminal decisions. For one such choice the unknowns are the stationary probabilities $\alpha_z(a)$ and $k(j\mid z,a,y)$, each in a finite simplex. Enumerating length-at-most-$n$ physical/label paths expresses their probabilities as polynomials with rational coefficients. Requiring zero probability of every premature stop, illegal action and non-stop at $n$ gives polynomial equalities. Each forbidden-path probability is a sum of nonnegative monomials. Expected loss is a polynomial on this closed subset of the compact product of simplexes. Thus a minimum exists when the set is nonempty. Quantifier elimination over real closed fields decides feasibility, gives its optimal algebraic value and an algebraic minimizing point; take the minimum over the finitely many enumerated choices\cite{bpr}. This procedure may be extremely expensive.

For each probability row of length $r$, round to integer numerators totaling $Q$ without introducing mass outside its original support. Taking floors and distributing the remaining units only within that support gives total variation at most $r/Q$. Coupling the action and label rows at each step gives mismatch probability at most $(A+M)/Q$. Crucially, support-preserving rounding cannot create any forbidden physical/label path: every such path has a zero raw factor or a zero original controller factor, and these zero factors remain zero. Hence legality and exact stopping hold with probability one for the rounded machine, not merely approximately. Proposition~\ref{prop:numerical} proves \eqref{eq:compiler_error}.

Store one action row per label and one next-label row for each $(z,a,y)$, together with stopping flags and decisions. This gives the displayed bit bound. To sample a denominator-$Q$ row, draw $\lceil\log_2Q\rceil$ fresh bits and reject values outside $\{0,\ldots,Q-1\}$. Acceptance probability is greater than one half, and cumulative sums find the selected outcome by scanning the row. Only an index, a sum and the sampled integer are needed in scratch space. The expected number of trials is less than two; row search uses $O((A+M)\log(Q+1))$ elementary bit operations per paired action/transition sample, in addition to the raw physical call. After $k$ rejections use a fixed member of the same support. Each fallback has probability at most $2^{-k}$, so two samples per call add at most $2n2^{-k}$ to the coupling defect. The bounded-trial variant additionally uses $O(\log(k+1))$ scratch bits for an erased trial counter. The number of sampling bit operations is then bounded by $O(n(k+A+M)\log(Q+1))$. The fallback never introduces a new support edge and therefore preserves the exact interface. A trial counter is scratch space erased before the next physical call, not a persistent random seed.
\end{proof}
The general continuous-instrument results do not inherit this finite rational decidability. Conversely, the finite theorem does not replace the main compact-carrier criterion by an enumeration of one special model.

\begin{corollary}[Attained acquisition, memory and error coordinates]\label{cor:four_resources}
Take a regular engine of Theorem~\ref{thm:regular}, or a blind engine of Theorem~\ref{thm:blind}. Before its $n\ge1$ calls, make one paid early call observing an independent fair bit through an erasure channel with erasure probability $2\varepsilon$, $0\le\varepsilon\le1/2$. An estimate of this bit must be written immediately to a write-only scoring port. Also prepare an independent offset taking $\pm\delta$ with equal probability, $0\le\delta\le1$, never reported to the predictor. Its squared prediction error is scored together with the engine's loss and the early bit error. One final paid audit is performed by the scoring authority after the predictor has halted and fixed all decisions, without releasing information beforehand. Thus $N=n+2$ counts physical acquisitions and the predictor halts after $N-1$ calls; it does not make a further transition for the audit.

The source-to-perfect-bit causal deficiency at the early completion deadline is exactly $\varepsilon$, uniformly over the two bit values. For the autonomous class, all $M\ge N$ satisfy, in the regular case,
\begin{equation}\label{eq:four_regular}
 \mathcal R(N,M,\delta,\varepsilon)-B_{N-2}^*
 \asymp (M-N+1)^{-2/d}+\delta^2+\varepsilon,
\end{equation}
and, in the blind case, the exact identity is
\begin{equation}\label{eq:four_blind}
 \mathcal R(N,M,\delta,\varepsilon)
 =q_{\lfloor(M-2)/(N-2)\rfloor}(\mu)+\delta^2+\varepsilon/2.
\end{equation}
The constants in \eqref{eq:four_regular} are those of the fixed regular class, independent of these budgets and errors. The calibration and deficiency terms are actual inaccessible alternatives and actual interface defects, not arbitrary permitted error tolerances.
\end{corollary}
\begin{proof}
The best bit prediction on erasure is $1/2$, giving squared loss $\varepsilon/2$; the offset has conditional mean zero and variance $\delta^2$. These variables are independent of the engine and scored additively under one product law, so their lower bounds add to its risk for any observer. The physical interface permits only the early action at the initial label and only engine actions afterwards. That initial label is therefore distinct from all engine labels. The write-only output is not readable and cannot function as a persistent register. Any extra bit-dependent randomization of the engine can be generated as fresh independent initial randomization; it supplies no information about the independent engine variable and cannot improve the minimum over initial engine programs. Conditional on that seed, the corresponding legal engine begins in one of at most $M-1$ labels. Taking the best such conditional program proves the engine lower with $M-1$. The matching upper uses one early label and the corresponding engine machine with $M-1$ labels; the bit estimate is emitted during the early transition. After the predictor halts and fixes its decisions, the scoring authority performs the final physical audit; this adds one charged acquisition but no predictor transition. Thus $N$ counts physical acquisitions, while the autonomous predictor itself halts after $N-1$ calls. This distinction is part of the stated task, not a free terminal signal during the engine. Applying the two engine laws gives the displayed formulas.

For the deficiency, the two erased-bit source report laws have TV distance $1-2\varepsilon$, whereas the perfect-bit laws have distance one. Any common parameter-independent postprocessing preserves the former upper bound. If its error from each perfect law were at most $e$, the triangle inequality would give $1\le2e+1-2\varepsilon$, hence $e\ge\varepsilon$. Reporting the observed bit and guessing fairly on erasure attains error $\varepsilon$ for both values. The imposed early deadline forbids using the later audit. Physical preparation, erasures, engine calls and the final scoring call are all counted; the calibration offset is not an unknown kernel that may be learned from hidden free samples.
\end{proof}

\section{Relation to earlier theories and scope}\label{sec:literature}
The static barycentric step belongs to the theory of comparison of experiments and convex order. Blackwell's comparison theorem\cite{blackwell} concerns the decision content of a parameter-independent garbling; Strassen's theorem\cite{strassen} identifies convex order with a martingale coupling. Lemma~\ref{lem:allocation} is a finite-target proof of this classical mechanism. Bayesian persuasion\cite{kamenica} optimizes information structures subject to Bayes plausibility; coarse information design\cite{lyu} relates conditional-mean cells to the curvature of their value function. Here the incoming law is fixed by the actually executed physical acquisition. One may only garble that law, and the next law depends on the retained information and chosen action. Neither Bayes plausibility alone nor a static posterior partition enforces these serial constraints.

Causal transport\cite{backhoff} imposes nonanticipation on couplings and has its own dynamic programming principle. The present restriction is different from nonanticipation alone: retained information must have prescribed finite cardinalities, and a reused internal label must invoke one and the same report kernel at every phase. Equation~\eqref{eq:simdual} imposes this latter constraint simultaneously, rather than selecting an arbitrary adapted coupling at each time. The theorem does not identify directed information, a channel rate, or a transport cost with persistent label cardinality. The compact finite-obstruction result is a topological separation statement; it is not an efficient optimization algorithm over causal transport plans.

Several approximation theories quantify finite-memory performance for known POMDPs. Kara and Y\"uksel\cite{kara} construct finite-window approximations under filter stability and finite observation/action alphabets. Demirci, Kara and Y\"uksel\cite{refined} refine such estimates using expected Wasserstein or pathwise total-variation stability. Saldi, Y\"uksel and Linder\cite{saldi} establish finite-model approximation for discounted partially observed control under continuity conditions. These approaches compare approximate policies or values with their control benchmarks. They do not impose the same internally counted exact-deadline interface used in \eqref{eq:regular_aut}. Our quantitative theorem instead combines an actual-law converse with explicitly allocated finite observer states, while the exact allocation theorem itself does not require their stability conditions. This distinction does not imply that the present theorem solves the discounted or average-cost problems treated in those works.

Grover and Dimitrakakis\cite{grover} adapt belief discretization to depth in POMDP planning. Li, Hammar and Bertsekas\cite{li} aggregate beliefs through features and conditional means and derive value-approximation bounds. Adaptive grid size, feature aggregation and conditional means are therefore not new ingredients here. The additional requirement is calibration under all incoming edges of the deployed observer, and for an internally timed observer the compatibility of the same transition programme across all cuts. The mass--curvature estimate and block allocation used in Section~\ref{sec:geometry} are inherited quantitative mechanisms\cite{prior}; they are not counted again as new results merely because they now follow an exact feasibility theorem.

The recent belief-metric covering analysis of Zhu and Lu\cite{zhu} addresses off-policy evaluation and coverage/sample-efficiency bounds for stable policies in offline POMDP learning. Its geometric coverage condition is not the hard cardinality of a causal observer's persistent state. Kim\cite{kim} studies finite-memory belief mismatch and policy-conditional performance along a common closed-loop execution. That information pattern differs both from optimizing over every legal controller and from charging a stationary machine's timing states. These comparisons are at the stated theorem objectives; neither paper is dismissed on the basis of its terminology or publication date. In particular, this article does not claim matching acquisition complexity for learning an unknown raw kernel from offline data.

Zero-delay source coding is another close but distinct theory. Wood, Linder and Y\"uksel\cite{wood} and Ghomi, Linder and Y\"uksel\cite{ghomi} study optimal stationary codes and finite-memory approximations for Markov sources under long-run distortion criteria. Stavrou, Charalambous, Charalambous and Loyka\cite{stavrou} develop nonanticipative rate-distortion bounds and causal filtering constructions, including Gaussian processes. An information-rate or communication-alphabet constraint does not bound every encoder/decoder internal state. Conversely, our observer need not transmit one symbol over a communication channel at every physical call. Thus their infinite-horizon and information-rate statements cannot be substituted for the finite-label cut converse here, and our results do not subsume theirs.

Positive quantum instruments have a standard operational basis\cite{davies}. Theorem~\ref{thm:quantum} uses a deliberately explicit finite-dimensional positive preparation and spanning instrument, so that quotient separation, report continuity and acquired density can be verified without commutativity. It is not a theorem about arbitrary quantum memory, operator domains, or physically unrestricted measurement design. Likewise, Proposition~\ref{prop:simulation} is the usual bounded-risk comparison mechanism\cite{lecam} with the missing causal and resource interfaces made explicit, rather than a new total-variation inequality. The effective finite compilation theorem uses classical real algebraic decision methods\cite{bpr}; it gives no new complexity bound for quantifier elimination.

The new proposed mathematical contribution is the exact finite-support causal allocation frontier, its phase-compatible dual, the compact global finite-obstruction consequence, and their combination with explicitly counted implementation and distinct matching regimes. The zero-loss criterion is independent of full task curvature; the block-stable law and blind-transport law exhibit when, and why, quantitative conclusions differ. Full predictive quotients of arbitrary Borel experiments, a uniformly effective theory of all stationary compact-carrier optimizers, unknown-kernel minimax learning, and a matched optimal region for programme size, scratch space and physical time are not established. Neither the fixed-horizon theory nor its bounded-stopping error estimate is an infinite-horizon theorem.

\appendix
\section{From raw experiments to a measurable predictive quotient}\label{sec:quotient}
We give a sufficient realization result at the raw level. It does not assert that every standard Borel predictive equivalence relation has a compact or even standard Borel quotient.

For each phase let $X_t$ be compact metrizable. Let
$K_t^a(x,dx'\,dy)$ be a positive normalized Borel physical kernel with standard Borel reports. Fix on each report space a compatible Polish topology and a countable bounded continuous measure-determining family containing constants and closed under rational linear combinations and products. Let $\mathcal V_t\subset C(X_t)$ be a separable closed linear subspace containing constants and the conditional expectations of a determining family of executable future tests. Assume it contains every pullback
\begin{equation}\label{eq:raw_pullback}
 x\longmapsto\int \varphi(y)f(x')K_t^a(x,dx'\,dy),
 \qquad f\in\mathcal V_{t+1},
\end{equation}
for legal $a$ and each function $\varphi$ in the chosen report family. In particular these pullbacks are continuous. Assume the declared future-test system is determined by these tests and their successive report-contingent continuations; equality on them must not discard a legal action or an observable resource coordinate. Phase-dependent legality is retained separately. The terminal task has continuous affine expectation in the resulting coordinates. These are verifiable closure conditions on the raw experiment, not assumptions of a Fisher metric or an available exact posterior register.

Choose a countable norm-dense family $f_{t,i}$ in the unit ball of $\mathcal V_t$, including constants, and define
\begin{equation}\label{eq:raw_quotient}
 Q_t:\mathcal P(X_t)\longrightarrow[-1,1]^{\mathbb N},\qquad
 Q_t(p)=(\int f_{t,i}\,dp)_{i\ge1},\qquad
 \mathsf K_t=Q_t(\mathcal P(X_t)).
\end{equation}
The dense family may be chosen from the rational span of determining executable tests; equality of its coordinates is therefore equality of their predictions, not an extra observable added to the experiment.

\begin{theorem}[Predictive realization and minimality]\label{thm:quotient}
Under the preceding raw hypotheses, $\mathsf K_t$ is compact metrizable and convex. Two physical distributions have the same $Q_t$ precisely when their declared executable future tests agree. The actual posterior of every legal finite-history strategy projects to a Borel $\mathsf K_t$-valued predictive state. Report and update kernels on $\mathsf K_t$ admit jointly Borel versions, independent of the representing distribution up to the relevant report-null sets, and satisfy \eqref{eq:affine_instrument}. Conditional mixing is realized by barycenters.

Every measurable statistic determining all these future tests factors the projected predictive state through a Borel map; that factor is unique on the statistic's attained values. Conversely, the projected state determines the tests. This is minimality for the fixed preparation/test interface, not simultaneous frequentist sufficiency for all inaccessible parameters. If the induced acquired-law kernels are weakly continuous, all the hypotheses concerning predictive realization in Theorem~\ref{thm:completion} hold.
\end{theorem}
\begin{proof}
Weak compactness of $\mathcal P(X_t)$ and continuity of the coordinates make $\mathsf K_t$ compact in the countable product; it is convex by linearity. Equality of dense coordinates gives equality on $\mathcal V_t$. The report pullbacks with $f=1$ determine the report law. With general $f$ they determine the finite signed measures
$B\mapsto\int_{X_t\times X_{t+1}\times B}f(x')p(dx)K_t^a(x,dx'\,dy)$.
The measure-determining report family and a monotone-class argument extend the identities to all Borel report sets. After disintegration, equal incoming coordinates therefore give equal outgoing conditional coordinates outside one report-null set, chosen common for the countable family. Induction over the finite continuation proves equality of the declared future tests, including report-contingent choices. The reverse implication holds because the determining tests are executable. Continuity of arbitrary report-contingent test expectations is not needed for this implication.

For measurability, a continuous map from a compact metric space onto a compact metric image has a Borel section here. An explicit construction is useful. Cover $\mathcal P(X_t)$ at level $m$ by finitely many closed balls of radius $2^{-m}$. Select the least ball whose intersection with the previously selected compact intersections still meets the fiber $Q_t^{-1}(q)$. The condition is Borel in $q$, since the image of each such compact intersection is closed. The nested nonempty compact intersections have diameters tending to zero, and their unique limiting point is a Borel selection $p(q)$ in the fiber.

Apply a jointly measurable conditional-probability construction to the raw kernel with input $(p(q),a)$. Such a construction follows, for example, by conditioning on successive finite partitions generating the report sigma field: conditional expectations of a countable dense family in $C(X_{t+1})$ converge by the martingale theorem, determine a probability measure on the compact state space almost everywhere, and have measurable limits. Assign an arbitrary fixed probability where the convergence or consistency tests fail. Project this posterior by $Q_{t+1}$ to define $F_t(q,a,y)$, and use the raw report marginal to define $J_t(q,a)$. The preceding signed-measure identities prove representative independence almost everywhere. The initial and all later history spaces are standard Borel, so this construction also gives Borel projected posteriors along actual histories.

For a probability $\nu$ on $\mathsf K_t$, mix its selected representatives: $p=\int p(q)\nu(dq)$. Then $Q_t(p)=\bar\nu$. Linearity of the unnormalized raw kernel in $p$, followed by the signed-measure identities, gives both formulas \eqref{eq:affine_instrument}. Equality on coordinate functions is enough: it identifies each conditional barycenter in $\mathsf K_{t+1}$ and hence every continuous affine function. This is mathematical realizability as a conditional physical mixture. It does not make re-preparation free or assert that every point of the compact carrier is reached under the chosen exploration. In a feasible allocation flow the induction of Lemma~\ref{lem:calibration} supplies the actual conditional mixtures under the deployed machine.

Finally, let a statistic $T(h)$ determine the tests, with Borel coordinate maps $b_i$ such that $Q_t(p_h)_i=b_i(T(h))$. The vector $b=(b_i)$ is Borel. Since $\mathsf K_t$ is closed in the product, replace $b$ by a fixed point outside the Borel set $b^{-1}(\mathsf K_t)$. This defines a Borel factor through $T$ agreeing with the projected state on every attained history where the conditional versions are specified. Its values there are forced by the separating coordinates. No claim that the image of an arbitrary history map is Borel is required. The remaining weak-continuity condition is additional; it is not a consequence of measurability alone.
\end{proof}
In the finite hidden-state model, the test span identifies the simplex of distributions. In the noncommuting instrument model, all future-test expectations are affine in the density matrix and the spanning effects identify that matrix; different classical decompositions of it are correctly identified. For capture of a compactly supported variable, bounded continuous report tests identify its probability law, and the carrier can be infinite-dimensional. These identifications explain why the same abstract theorem covers the three settings without assigning an ambient Fisher or covariance geometry at the outset.

\begin{thebibliography}{99}
\bibitem{backhoff} J. Backhoff, M. Beiglb\"ock, Y. Lin and A. Zalashko, \emph{Causal transport in discrete time and applications}, SIAM J. Optim. \textbf{27} (2017), 2528--2562, doi:10.1137/16M1080197; arXiv:1606.04062v2.
\bibitem{bpr} S. Basu, R. Pollack and M.-F. Roy, \emph{Algorithms in Real Algebraic Geometry}, second ed., Algorithms and Computation in Mathematics 10, Springer, 2006, doi:10.1007/3-540-33099-2.
\bibitem{blackwell} D. Blackwell, \emph{Equivalent comparisons of experiments}, Ann. Math. Statist. \textbf{24} (1953), 265--272, doi:10.1214/aoms/1177729032.
\bibitem{davies} E. B. Davies and J. T. Lewis, \emph{An operational approach to quantum probability}, Comm. Math. Phys. \textbf{17} (1970), 239--260, doi:10.1007/BF01647093.
\bibitem{refined} Y. E. Demirci, A. D. Kara and S. Y\"uksel, \emph{Refined bounds on near optimality finite window policies in POMDPs and their reinforcement learning}, arXiv:2409.04351v1 (2024).
\bibitem{ghomi} M. Ghomi, T. Linder and S. Y\"uksel, \emph{Zero-delay lossy coding of linear vector Markov sources: optimality of stationary codes and near optimality of finite memory codes}, IEEE Trans. Inform. Theory \textbf{68} (2022), 3474--3488, doi:10.1109/TIT.2021.3138769; arXiv:2103.10810v3.
\bibitem{grover} D. Grover and C. Dimitrakakis, \emph{Adaptive belief discretization for POMDP planning}, arXiv:2104.07276v1 (2021).
\bibitem{kamenica} E. Kamenica and M. Gentzkow, \emph{Bayesian persuasion}, Amer. Econ. Rev. \textbf{101} (2011), 2590--2615, doi:10.1257/aer.101.6.2590.
\bibitem{kara} A. D. Kara and S. Y\"uksel, \emph{Near optimality of finite memory feedback policies in partially observed Markov decision processes}, J. Mach. Learn. Res. \textbf{23} (2022), no. 11, 1--46.
\bibitem{kim} M. Kim, \emph{Finite memory belief approximation for optimal control in partially observable Markov decision processes}, arXiv:2601.03132v1 (2026).
\bibitem{lecam} L. Le Cam, \emph{Asymptotic Methods in Statistical Decision Theory}, Springer Series in Statistics, Springer, 1986.
\bibitem{li} Y. Li, K. Hammar and D. Bertsekas, \emph{Feature-based belief aggregation for partially observable Markov decision problems}, arXiv:2507.04646v1 (2025).
\bibitem{lyu} Q. Lyu, W. Suen and Y. Zhang, \emph{Coarse information design}, arXiv:2305.18020v5 (2025).
\bibitem{prior} Q. Qi, \emph{General Theta Foundations I: Acquired Geometry and Causal Resource Transfer}, complete earlier mass--curvature article, archival companion V, 2026. The posterior-pooling, adaptive-completion and earlier technical articles are preserved separately as archival companions U, T and S; none is required to supply an unproved premise of the present article.
\bibitem{saldi} N. Saldi, S. Y\"uksel and T. Linder, \emph{Finite model approximations for partially observed Markov decision processes with discounted cost}, IEEE Trans. Automat. Control \textbf{65} (2020); preprint version arXiv:1710.07009v1 (2017).
\bibitem{stavrou} P. A. Stavrou, T. Charalambous, C. D. Charalambous and S. Loyka, \emph{Optimal estimation via nonanticipative rate distortion function and applications to time-varying Gauss--Markov processes}, SIAM J. Control Optim. \textbf{56} (2018), 3731--3765, doi:10.1137/17M1116349.
\bibitem{strassen} V. Strassen, \emph{The existence of probability measures with given marginals}, Ann. Math. Statist. \textbf{36} (1965), 423--439, doi:10.1214/aoms/1177700153.
\bibitem{wood} R. G. Wood, T. Linder and S. Y\"uksel, \emph{Optimal zero delay coding of Markov sources: stationary and finite memory codes}, IEEE Trans. Inform. Theory \textbf{63} (2017), 5968--5980; arXiv:1606.09135v2.
\bibitem{zhu} Y. Zhu and Y. Lu, \emph{A covering framework for offline POMDPs learning using belief space metric}, Proc. 29th International Conference on Artificial Intelligence and Statistics, Proc. Mach. Learn. Res. \textbf{300} (2026), 1423--1431.
\end{thebibliography}

\end{document}
